\documentclass[12pt,a4paper]{article}

\usepackage[french]{babel}
\usepackage[utf8]{inputenc}
\usepackage[T1]{fontenc}
\usepackage{geometry}
\usepackage{booktabs}
\usepackage{longtable}
\usepackage{array}
\usepackage{hyperref}
\usepackage{setspace}
\usepackage{parskip}

% Allow URLs to break across lines in bibliography
\def\UrlBreaks{\do\/\do-\do\_\do.\do:}

\geometry{
  top=2.5cm,
  bottom=2.5cm,
  left=3cm,
  right=2.5cm
}

\setlength{\parskip}{6pt}
\setlength{\parindent}{0pt}

\hypersetup{
  colorlinks=true,
  linkcolor=black,
  urlcolor=blue,
  citecolor=black
}

\begin{document}

% ─── PAGE DE TITRE ───────────────────────────────────────────────────────────
\begin{titlepage}
  \centering
  \vspace*{2cm}
  {\Large\bfseries EMLV\\[0.4em]MSc Digital Business Analytics\par}
  \vspace{2cm}
  {\LARGE\bfseries Revue de littérature\par}
  \vspace{1cm}
  {\Large Comment l'intelligence artificielle influence-t-elle le\\[0.3em]
  bien-être du consommateur~?\par}
  \vspace{0.5cm}
  {\large Entre gains d'efficacité et risques liés aux données\par}
  \vspace{2cm}
  {\large Othmane Tazouti\par}
  \vspace{1cm}
  {\large Année académique 2025--2026\par}
\end{titlepage}

% ─── TABLE DES MATIÈRES ──────────────────────────────────────────────────────
\tableofcontents
\newpage

% ─── 1. INTRODUCTION ─────────────────────────────────────────────────────────
\section{Introduction}

L'intelligence artificielle (IA) bouleverse complètement les relations entre entreprises et
individus, qu'il s'agisse de consommateurs, d'employés ou de citoyens. Aujourd'hui, avec
la généralisation des systèmes algorithmiques dans les entreprises pour optimiser leurs
opérations, le bien-être des parties prenantes devient un enjeu capital. Cette revue de
littérature s'attache justement à cette problématique : Comment l'IA influence-t-elle le
bien-être du consommateur, entre gains d'efficacité et risques liés aux données~?

Pour y répondre, on s'appuie sur 23~articles scientifiques récents (2022--2025) qui étudient
l'IA dans la gestion des ressources humaines~(GRH), le comportement organisationnel
et l'expérience des parties prenantes. Même si ces travaux se concentrent surtout sur le
domaine organisationnel, ils révèlent des mécanismes transposables au consommateur~:
personnalisation algorithmique, analyse prédictive, automatisation des services, surveillance
des données et enjeux éthiques. Une étude bibliométrique pionnière de Palos-Sánchez
et al.~(2022) a établi les jalons de ce champ, montrant que la recherche sur l'IA
appliquée à la GRH est en croissance constante, avec une concentration encore forte sur
le recrutement et la sélection. La bibliométrie récente de Soulami et al.~(2024) confirme
cette dynamique et révèle quatre grands clusters thématiques -- éthique, autonomie au
travail, stress de l'employé et santé mentale -- illustrant la manière dont l'IA façonne
profondément l'environnement professionnel.

Cette revue se construit autour de cinq axes thématiques. D'abord, on s'intéresse aux gains
d'efficacité grâce à l'IA pour optimiser les process et améliorer l'expérience utilisateur.
Ensuite, on analyse les risques liés aux données, dont la protection de la vie privée et les
biais algorithmiques. Troisième axe, l'engagement et le bien-être individuel face à l'IA.
Quatrième volet, les craintes de substitution et l'anxiété technologique. Cinquième, les
cadres pour une IA centrée sur l'humain. Enfin, une synthèse critique et quelques pistes
pour des recherches futures.

% ─── 2. GAINS D'EFFICACITÉ ───────────────────────────────────────────────────
\section{Gains d'efficacité et optimisation par l'IA}

\subsection{Automatisation des processus et réduction des délais}

Parmi les bénéfices les plus étudiés de l'IA, on trouve l'automatisation des tâches répétitives
et la chute spectaculaire des délais opérationnels. Kayusi et al.~(2025), à travers leur enquête
quantitative auprès de 316~professionnels RH, montrent que l'analytique RH basée sur
l'IA a réduit le temps de recrutement de 51,1~\% et augmenté la précision des décisions de
50,8~\%. Ces chiffres rejoignent ceux de Chheda et al.~(2025), qui, après l'analyse de
1~000~dossiers d'employés, constatent des gains notables en efficacité pour le recrutement et
la rétention grâce aux algorithmes prédictifs.

Sur l'expérience des employés, Mahmood et al.~(2025) rapportent que 75~\% des organisations
consultées utilisent l'IA pour le recrutement et 85~\% ont automatisé la gestion de la paie.
L'intégration technologique prend clairement une grande ampleur. Pandit~(2025) et
Biswas et al.~(2025) confirment que l'IA accélère l'accès aux talents, réduit les erreurs
humaines et optimise l'allocation des ressources.

Dans une étude analytique fondée sur une revue de littérature couvrant la période 2019--2025,
Almuhanna~(2025) montre que l'IA redéfinit en profondeur les fonctions RH~-- du recrutement
à la gestion des talents, en passant par l'évaluation des performances et la rémunération~--
en améliorant la précision décisionnelle et l'équité organisationnelle. L'auteure souligne
toutefois que ces bénéfices s'accompagnent de défis liés à l'éthique, à la transparence
algorithmique et à la gouvernance des données. Dans le contexte saoudien, elle observe
une disposition avancée à adopter des solutions RH pilotées par l'IA, conformément à la
Vision 2030, tout en insistant sur la nécessité de développer davantage les compétences
numériques.

Chez le consommateur, ces gains se traduisent par un service plus rapide, des recommandations
personnalisées et une prise en charge des demandes accélérée, ce qui améliore
nettement le bien-être ressenti.

\subsection{Personnalisation et analyse prédictive}

L'IA permet une personnalisation jamais vue des interactions. Singh~(2025) explique que les
systèmes IA analysent des tonnes de données comportementales pour adapter les parcours
de formation et développement à chaque profil. Kessavane~(2025) souligne l'utilisation du
machine learning pour des évaluations en continu, apportant un feedback en temps réel et
sur mesure.

Rathore et Rathore~(2024) détaillent comment le traitement du langage naturel~(NLP) et
le machine learning servent à prédire le turnover ou analyser le ressenti, des outils d'abord
déployés dans les RH mais qui trouvent des équivalents directs pour le client~: chatbots,
moteurs de recommandation, systèmes de tarification dynamique reposent sur ces mêmes
bases.

Sharma et Gupta~(2025), après une étude mixte auprès de 320~RH indiens et une modélisation
PLS-SEM, montrent que l'IA améliore la décision et la satisfaction quand les systèmes
paraissent transparents et équitables. Un point clé pour l'adoption par le consommateur.

Dans le domaine de la formation, Wagh~(2025) montre que l'IA et l'automatisation
révolutionnent le développement des compétences en entreprise, en offrant des plateformes
d'apprentissage adaptatif qui personnalisent le contenu selon le style d'apprentissage,
les lacunes et les aspirations de chaque salarié. Des tendances émergentes comme le coaching
virtuel et la gestion de la performance pilotée par l'IA témoignent d'une montée en
sophistication de ces outils. Benabou et Touhami~(2025a), dans une analyse bibliométrique
et systématique de 77~articles, corroborent ces résultats en identifiant le rôle transformateur
de l'IA dans la GRH -- notamment en formation et développement -- comme l'un des trois
grands axes de recherche du domaine.

% ─── 3. RISQUES LIÉS AUX DONNÉES ────────────────────────────────────────────
\section{Risques liés aux données et enjeux éthiques}

\subsection{Protection de la vie privée et surveillance algorithmique}

L'IA fait gagner du temps et de l'efficacité, mais elle soulève aussi de gros soucis en matière
de données personnelles. Ateeq et al.~(2025) étudient l'impact de l'IA sur les contrats
psychologiques en entreprise, mettant en garde contre la surveillance excessive. Ils montrent
que la collecte permanente de données comportementales peut miner la confiance et nuire
au bien-être psychologique.

Pandit~(2025) pointe la conformité au RGPD et les questions de confidentialité comme des
défis majeurs. Biswas et al.~(2025) insistent sur le fait que le manque de transparence
freine l'acceptation des nouvelles technologies. Sundari et al.~(2025), dans une revue
systématique PRISMA de 78~articles, concluent que la protection des données et la
gouvernance sont des conditions indispensables pour un usage responsable de l'IA.

Benabou et Touhami~(2025b), dans leur revue systématique PRISMA de 141~études publiées
entre 2019 et 2025, confirment que la sécurité des données et la prise de décision éthique
constituent des défis persistants. Leurs travaux montrent que l'IA offre des avantages
significatifs en matière d'efficacité et d'engagement, mais que des défis subsistent quant
aux décisions éthiques et au maintien de l'interaction humaine dans les processus RH.
Cette double nature -- opportunités et risques -- rejoint exactement la tension identifiée
dans la présente revue entre gains d'efficacité et protection des données.

Chez le consommateur, cela revient à arbitrer entre le bénéfice d'une offre personnalisée
(qui nécessite des données) et la préservation de son intimité. Ce fameux \emph{privacy
paradox}, où on affirme tenir à sa vie privée tout en acceptant de partager ses infos, reste
un champ de recherche clé.

\subsection{Biais algorithmiques et équité}

Les biais algorithmiques, c'est une menace directe pour l'équité -- et pour le bien-être
des personnes discriminées. Bankins et al.~(2024), dans une revue multiniveaux sur
68~études empiriques, identifient les biais dans les systèmes d'aide à la décision comme un
des dangers principaux de l'IA en entreprise. Ces biais sont présents à tous les niveaux
-- individuel, collectif, organisationnel -- et reflètent souvent les préjugés des données
d'apprentissage.

Croitoru et al.~(2025), avec le modèle TAM et 150~répondants, démontrent que la perception
d'équité est un facteur clé d'acceptation de l'IA. Kayusi et al.~(2025) recommandent des
audits réguliers des algorithmes pour détecter et corriger ces biais. Pour le consommateur,
cela peut se traduire en tarification discriminatoire, recommandations biaisées ou exclusion
de certains groupes.

% ─── 4. ENGAGEMENT, SATISFACTION ET BIEN-ÊTRE ───────────────────────────────
\section{Engagement, satisfaction et bien-être face à l'IA}

\subsection{Mécanismes d'engagement et conditions psychologiques}

Tortorella et al.~(2025) analysent l'impact de l'IA sur l'engagement des employés dans des
organisations lean, à partir du modèle de Kahn~(1990) avec trois conditions psychologiques~:
sens, sécurité, disponibilité. Les résultats montrent que l'IA renforce l'engagement -- en
libérant du temps pour des tâches à haute valeur -- mais peut aussi fragiliser celui-ci si
les individus perdent le contrôle ou le sens dans leur travail.

Qian et al.~(2025), en s'appuyant sur la théorie de la conservation des ressources~(COR) et
350~questionnaires, indiquent que l'usage de l'IA au travail stimule l'innovation, mais que
cet effet dépend du soutien organisationnel et du sentiment de justice. Pour le consommateur
aussi, l'engagement envers une marque ou un service grandit si l'IA améliore l'expérience
sans nuire au sentiment de contrôle.

Soulami, Benchekroun et Galiulina~(2024), dans une revue bibliométrique et systématique
publiée dans \emph{Frontiers in Artificial Intelligence}, analysent 92~articles via
VOSviewer et approfondissent l'examen de 25~études sélectionnées selon la méthode
PRISMA. Leur analyse co-occurrentielle dégage quatre grands clusters thématiques au croisement
IA--bien-être au travail~: l'éthique, l'autonomie au travail, le stress de l'employé et la santé
mentale. Ces quatre dimensions illustrent la complexité des dynamiques créées par l'IA
dans l'environnement professionnel. Les auteurs concluent que les organisations doivent
mettre en place des stratégies de soutien pour atténuer les effets négatifs potentiels de l'IA
sur le bien-être des employés, tout en tirant parti de ses bénéfices pour renforcer
l'autonomie et la satisfaction, et pour promouvoir l'innovation permise par l'IA grâce à la
créativité et au sentiment d'auto-efficacité des collaborateurs. Ces conclusions, obtenues
dans un contexte organisationnel, s'appliquent directement aux consommateurs~: les
services algorithmiques bien conçus peuvent stimuler l'autonomie et la satisfaction, tandis
que ceux qui sont opaques ou contraignants génèrent stress et méfiance.

Liu et Chen~(2025), dans une étude empirique menée auprès de 202~répondants à
l'Université de Nanjing Tech~(Chine), s'appuient sur la théorie personne-environnement
(\emph{person-environment fit}) pour explorer l'impact de l'usage de l'IA sur l'engagement
professionnel (\emph{career commitment}). Ils montrent que l'usage de l'IA a un effet
positif sur l'engagement professionnel, et que ce lien est médiatisé par le \emph{job crafting}~:
les employés qui adaptent proactivement leurs tâches en réponse aux outils IA renforcent
leur identité professionnelle et leur engagement à long terme. Ce résultat souligne que
l'effet positif de l'IA sur le bien-être n'est pas automatique~; il passe par la capacité
des individus à se réapproprier leur travail, un mécanisme transposable à la relation
consommateur--service algorithmique.

\subsection{Satisfaction et perception de la qualité de service}

Kayusi et al.~(2025) observent une hausse de 51,3~\% de la satisfaction des parties prenantes
après l'implémentation d'analytique RH basée sur l'IA. Mahmood et al.~(2025) confirment
que les technologies RH améliorent vraiment l'expérience employé quand elles sont bien
intégrées. Sharma et Gupta~(2025) ajoutent que la satisfaction dépend de la confiance
dans le système et du sentiment de transparence.

Sharma et al.~(2025), après une analyse bibliométrique sur plus de 11~291~articles, notent
que l'engagement et la satisfaction sont les thèmes de recherche les plus dynamiques à
l'intersection IA--comportement organisationnel. Cette tendance montre qu'on accorde de
plus en plus d'importance au facteur humain dans le déploiement de l'IA -- une réalité
qui vaut aussi pour la relation entre consommateur et services algorithmiques.

% ─── 5. ANXIÉTÉ TECHNOLOGIQUE ────────────────────────────────────────────────
\section{Anxiété technologique et craintes de substitution}

\subsection{Peur du remplacement et perte de contrôle}

L'anxiété face à l'IA impacte sérieusement le bien-être. Qian et al.~(2025), grâce à la
théorie COR, montrent que le sentiment de menace sur ses ressources (compétences, emploi,
autonomie) provoque du stress et réduit le bien-être. Bankins et al.~(2024) remarquent
que la peur d'être remplacé par l'IA affecte la motivation, l'identité professionnelle et le
contrat psychologique entre individu et organisation.

Ateeq et al.~(2025) détaillent comment l'IA modifie les contrats psychologiques en
entreprise. Ils observent un phénomène de «~rupture de contrat psychologique~» quand
l'IA est déployée sans explication claire, sans formation ou sans recours. Cela se transpose
facilement à la relation client--entreprise~: lancer un chatbot sans alternative humaine, ou
une décision opaque, peut être vécu comme une rupture d'engagement de service.

Liu et Chen~(2025) affinent cette analyse en introduisant la notion d'\emph{AI awareness}~--
la perception par l'individu de la menace que représente l'IA pour son développement de
carrière. Leurs résultats montrent que cette conscience joue un rôle modérateur négatif~:
un niveau élevé d'AI awareness réduit la propension au job crafting et atténue ainsi
l'effet positif de l'usage de l'IA sur l'engagement professionnel. Autrement dit, plus un
individu perçoit l'IA comme une menace, moins il est enclin à adapter proactivement son
travail, et moins il tire de bénéfices de l'adoption de ces outils. Ce mécanisme vaut
également pour le consommateur~: la méfiance envers les systèmes algorithmiques peut
freiner l'adoption et réduire les gains de bien-être qui en découlent.

\subsection{Nécessité de montée en compétences}

Face aux évolutions induites par l'IA, tout le monde s'accorde à dire qu'il faut investir
dans la formation et le développement des compétences. Sundari et al.~(2025) préconisent
un cadre IA centré sur l'humain avec des programmes de requalification. Singh~(2025) et
Biswas et al.~(2025) insistent sur la préparation des individus à collaborer avec les systèmes
intelligents plutôt que de les subir.

Croitoru et al.~(2025), avec le modèle TAM, montrent que l'utilité perçue et la facilité
d'utilisation sont des antécédents majeurs à l'acceptation de l'IA. Résultat~: la littératie
numérique du consommateur devient un facteur clé pour l'impact de l'IA sur son bien-être~;
un client informé et compétent profite mieux des services algorithmiques et protège ses
intérêts.

Dans ce même esprit, Wagh~(2025) préconise d'investir massivement dans les plateformes
d'apprentissage pilotées par l'IA, lesquelles permettent un suivi des résultats de formation
et une prévision de la progression de carrière. Il souligne par ailleurs l'émergence d'un
nouveau rôle pour les professionnels RH~: celui de «~facilitateur de la transformation numérique~»,
compétent en matière d'IA et de littératie des données. Ce rôle émergent a son équivalent
côté consommateur~: le «~consommateur averti~» doit maîtriser les codes des services
algorithmiques pour en tirer le meilleur parti tout en protégeant ses données.

% ─── 6. VERS UNE IA CENTRÉE SUR L'HUMAIN ────────────────────────────────────
\section{Vers une IA centrée sur l'humain}

\subsection{Cadres théoriques et gouvernance éthique}

Plusieurs auteurs proposent des cadres pour un déploiement éthique et centré sur l'humain
de l'IA. Sundari et al.~(2025) développent un cadre intégratif autour de quatre piliers~:
transparence, équité des décisions, protection des données et autonomie individuelle.

Bankins et al.~(2024) plaident pour une approche multiniveaux, car l'IA touche l'individu,
le groupe et l'organisation -- et les interventions doivent suivre.

Ateeq et al.~(2025) proposent un modèle de gouvernance éthique basé sur la préservation
du contrat psychologique. Ils recommandent la création de comités éthiques, des procédures
de recours et des mécanismes d'explicabilité des décisions IA. Tortorella et al.~(2025)
ajoutent que l'IA dans les organisations lean doit garder l'amélioration continue et le
respect des personnes.

Benabou et Touhami~(2025b), dans leur revue systématique PRISMA de 141~études, identifient
la collaboration homme--IA comme l'un des trois grands thèmes structurant la recherche
contemporaine sur l'IA en GRH. Ils recommandent d'augmenter -- plutôt que de remplacer~--
les rôles humains par des outils IA, et insistent sur la nécessité de cadres éthiques pour
guider l'intégration stratégique de ces technologies. Benabou et Touhami~(2025a) arrivent
aux mêmes conclusions dans leur analyse bibliométrique de 77~articles~: la collaboration
homme--IA émerge comme un paradigme structurant, et les organisations doivent naviguer
avec discernement entre les opportunités et les défis que pose l'IA. Ces deux contributions
complémentaires des mêmes auteurs, publiées dans des revues distinctes, renforcent
la robustesse de ce constat.

\subsection{Équilibre entre automatisation et intervention humaine}

La question de l'équilibre automation--humain revient dans tout le corpus. Sharma et
Gupta~(2025) montrent que les systèmes hybrides, associant IA et jugement humain,
donnent de meilleurs résultats que les systèmes purement automatisés. Kessavane~(2025)
recommande que les évaluations IA soient toujours validées par un décideur humain.

Pandit~(2025) et Singh~(2025) insistent sur la nécessité de garder l'humain «~dans la
boucle~» (\emph{human-in-the-loop}) pour les décisions sensibles. Pour le consommateur,
ce principe implique le droit de contester une décision IA, d'obtenir un interlocuteur
humain et de connaître les critères de personnalisation. L'article~22 du RGPD, qui protège
contre la décision entièrement automatisée, formalise cette exigence sur le plan réglementaire.

% ─── 7. SYNTHÈSE ET DISCUSSION ───────────────────────────────────────────────
\section{Synthèse et discussion critique}

\subsection{Tableau récapitulatif}

\begin{longtable}{>{\bfseries\raggedright}p{3.4cm} >{\raggedright}p{6.2cm} >{\raggedright\arraybackslash}p{4.6cm}}
\toprule
Thème & Principaux résultats & Sources clés \\
\midrule
\endhead
\bottomrule
\endfoot

Gains d'efficacité
  & Réduction du temps de recrutement (--51,1\,\%), précision (+50,8\,\%),
    automatisation (85\,\%~paie)
  & Kayusi et al. (2025), Mahmood et al. (2025), Chheda et al. (2025) \\[6pt]

Personnalisation
  & Analyse prédictive, NLP, feedback en temps réel, adaptation des
    parcours ; apprentissage adaptatif et coaching virtuel
  & Rathore \& Rathore~(2024), Singh~(2025), Kessavane~(2025),
    Wagh~(2025) \\[6pt]

Vie privée et données
  & Risques RGPD, surveillance, perte de confiance, \emph{privacy
    paradox}~; défis éthiques et de sécurité des données
  & Ateeq et al.~(2025), Pandit~(2025), Sundari et al.~(2025),
    Benabou \& Touhami~(2025b) \\[6pt]

Biais algorithmiques
  & Discrimination, audits nécessaires, perception d'équité
  & Bankins et al.~(2024), Croitoru et al.~(2025),
    Kayusi et al.~(2025) \\[6pt]

Engagement
  & Conditions de Kahn, théorie~COR, double effet de l'IA~;
    clusters éthique / autonomie / stress / santé mentale
  & Tortorella et al.~(2025), Qian et al.~(2025),
    Soulami et al.~(2024) \\[6pt]

Engagement professionnel
  & Usage IA → \emph{job crafting} → engagement~; \emph{AI awareness}
    modère négativement
  & Liu \& Chen~(2025) \\[6pt]

Satisfaction
  & +51,3~\% de satisfaction, poids de la confiance et de la transparence
  & Kayusi et al.~(2025), Sharma \& Gupta~(2025) \\[6pt]

Anxiété
  & Peur de substitution, rupture de contrat psychologique,
    \emph{AI awareness} comme frein
  & Qian et al.~(2025), Ateeq et al.~(2025), Bankins et al.~(2024),
    Liu \& Chen~(2025) \\[6pt]

IA centrée sur l'humain
  & Cadres éthiques, gouvernance, \emph{human-in-the-loop},
    collaboration homme--IA
  & Sundari et al.~(2025), Ateeq et al.~(2025), Sharma \& Gupta~(2025),
    Benabou \& Touhami~(2025a, 2025b) \\[6pt]

Formation \& transfor-\\mation RH
  & IA redéfinit les fonctions RH ; montée en compétences
    indispensable ; rôle émergent du facilitateur numérique
  & Almuhanna~(2025), Wagh~(2025),
    Benabou \& Touhami~(2025a) \\

\end{longtable}

\subsection{Discussion}

Cette analyse croisée des 23~articles fait ressortir une tension majeure au cœur de la
relation IA--bien-être~: les mêmes technologies qui optimisent et personnalisent sont aussi
celles qui menacent la vie privée et provoquent de l'anxiété. Voilà toute la problématique
«~entre gains d'efficacité et risques liés aux données~».

Quelques constats se dégagent. D'abord, les gains d'efficacité sont bien quantifiés~: les
études de Kayusi et al.~(2025), Chheda et al.~(2025) et Mahmood et al.~(2025) apportent
des données solides sur l'amélioration des process. Ensuite, les risques éthiques sont moins
souvent chiffrés mais systématiquement repérés~: chaque article cite au moins un problème
lié aux données. Enfin, la satisfaction et l'acceptation dépendent de la transparence, de
l'équité et du maintien du contrôle humain -- trois leviers pour les organisations qui
veulent maximiser le bien-être de leurs parties prenantes.

Les études de Soulami et al.~(2024) et Liu et Chen~(2025) apportent un éclairage
complémentaire~: le bien-être n'est pas un résultat passif de l'IA, mais une construction
active qui repose sur les stratégies d'adaptation (job crafting, montée en compétences) et
sur la qualité de l'environnement organisationnel (autonomie, soutien, transparence).
Benabou et Touhami~(2025a, 2025b) et Almuhanna~(2025) convergent vers le même
constat~: l'IA augmente plutôt qu'elle ne remplace l'humain, à condition que les
organisations investissent dans des cadres éthiques et des stratégies de collaboration
homme--machine.

Du point de vue consommateur, les résultats RH sont transférables~: mêmes mécanismes
de personnalisation, de prédiction et d'automatisation~; mêmes risques de biais, de
surveillance et de perte de contrôle~; et les cadres éthiques proposés (transparence,
explicabilité, recours humain) valent aussi bien pour les services B2C que pour les RH
internes.

\subsection{Limites de la littérature existante}

Le corpus est riche, mais il y a quelques limites. D'abord, la plupart des études regardent
surtout la GRH, ce qui réduit l'application directe au domaine de la consommation. Ensuite,
les études quantitatives (Kayusi et al.~2025~; Sharma \& Gupta~2025~; Croitoru et al.~2025)
reposent souvent sur des échantillons de convenance dans certains pays (Inde, Roumanie,
Chine), ce qui pose des questions de validité externe. Enfin, les études longitudinales
manquent~: la plupart des travaux donnent des images à l'instant~$t$ sans suivre l'évolution
des perceptions.

Il convient également de noter que les revues bibliométriques et systématiques (Soulami
et al.~2024~; Benabou \& Touhami 2025a, 2025b~; Sharma et al.~2025) offrent une vue
macroscopique utile mais doivent être complétées par des études empiriques primaires,
notamment sur les populations de consommateurs. L'article fondateur de Palos-Sánchez
et al.~(2022), bien qu'antérieur à la période 2024--2025, reste pertinent pour situer
historiquement la trajectoire de recherche sur l'IA en GRH~: il documente la montée en
puissance du champ depuis 2018 et sa concentration initiale sur le recrutement, qui s'est
depuis étendue à l'ensemble des fonctions RH et au bien-être des parties prenantes.

% ─── RÉFÉRENCES ──────────────────────────────────────────────────────────────
\newpage
\section*{Références bibliographiques}
\addcontentsline{toc}{section}{Références bibliographiques}

\sloppy
\begin{enumerate}

\item Ateeq, K., Oswal, N., Jawabri, A., Masaeid, T.~A., Alquqa, E.~K., Basha, S.~E.,
  \& Alami, R. (2025). The transformative impact of artificial intelligence~(AI) on
  organisational behaviour~(OB)~: A study of employee engagement, performance, and
  ethical implications. \textit{Journal of Posthumanism}, \textit{5}(4), 1245--1256.

\item Bankins, S., et al. (2024). A multilevel review of artificial intelligence in
  organizations~: Implications for organizational behavior. \textit{Journal of
  Organizational Behavior}, \textit{45}, 159--182.

\item Biswas, O., Sharma, P., Das, S., \& Bharti, L. (2025). A study on the impact of
  artificial intelligence on human resource management. \textit{International Journal of
  Scientific Research in Engineering and Management~(IJSREM)}, \textit{9}(4).

\item Chheda, K., Thakur, U., J., R., Prasad Das, G., Parashar, M.~K., \& Reddy, K.
  (2025). AI-powered human resource management for enhancing employee recruitment
  efficiency and talent retention in organizations. \textit{Management (Montevideo)},
  \textit{3}, 165.

\item Croitoru, G., et al. (2025). Artificial intelligence in human resource management
  and its impact on employee motivation. \textit{Annals of Dunarea de Jos University,
  Fascicle~I}.

\item Kayusi, S., et al. (2025). AI-driven HR analytics~: Enhancing workforce
  decision-making through predictive insights. \textit{REMUVAC}, \textit{214},
  558--582.

\item Kessavane, I. (2025). AI and machine learning in human resource data analytics for
  performance management. \textit{International Journal on Science and Technology
  (IJSAT)}, article n\textsuperscript{o}~3729.

\item Mahmood, A., et al. (2025). HR tech and employee experience~: Exploring the role
  of AI, automation, and people analytics in modern workplaces. \textit{Indus Journal
  of Social Sciences}, \textit{1}(1--3).

\item Pandit, V. (2025). The role of artificial intelligence in transforming human
  resource management. \textit{International Journal of Scientific Research in
  Engineering and Management (IJSREM)}, \textit{9}(2).

\item Qian, X., et al. (2025). How does AI usage in the workplace affect employees'
  innovative behavior~? \textit{Behavioral Sciences}, \textit{15}, 465.

\item Rathore, S., \& Rathore, A. (2024). Machine learning applications in human
  resource management. \textit{International Journal of Global Academic Scientific
  Research~(IJGASR)}, \textit{3}(2), 77.

\item Sharma, A., \& Gupta, S. (2025). Artificial intelligence in HRM~: A mixed-methods
  study. \textit{Indian Journal of Engineering Sciences}, \textit{11}(24s), 259.

\item Sharma, M., et al. (2025). Artificial intelligence and employee engagement~: A
  bibliometric analysis. \textit{Information~(MDPI)}, \textit{16}, 555.

\item Singh, P. (2025). A comprehensive review of artificial intelligence in human
  resource management. \textit{Brixton Scholarly Review}, 76--95.

\item Sundari, T., et al. (2025). A systematic literature review on AI in HRM~: Towards
  a human-centered framework. \textit{International Research Journal of Advanced
  Engineering and Management~(IRJAEM)}, \textit{3}(10), 3010--3022.

\item Tortorella, G., et al. (2025). How does artificial intelligence impact employee
  engagement in lean organisations~? \textit{International Journal of Production
  Research}.

  % ── Nouvelles références ──────────────────────────────────────────────────

\item Almuhanna, N.~A. (2025). The role of artificial intelligence in transforming
  human resource management processes~: An analytical study. \textit{Academic Journal
  of Research and Scientific Publishing~(AJRSP)}, \textit{7}(79).

\item Benabou, A., \& Touhami, F. (2025a). Empowering human resource management through
  artificial intelligence~: A systematic literature review and bibliometric analysis.
  \textit{International Journal of Production Management and Engineering~(IJPME)},
  \textit{13}(1), 59--76.
  \url{https://doi.org/10.4995/ijpme.2025.21900}

\item Benabou, A., \& Touhami, F. (2025b). Artificial intelligence in human resource
  management~: A PRISMA-based systematic review. \textit{Acta Informatica Pragensia},
  \textit{14}(3), 489--505.
  \url{https://doi.org/10.18267/j.aip.264}

\item Liu, X., \& Chen, Y. (2025). The impact of artificial intelligence usage on
  employee career commitment~: The moderating role of artificial intelligence awareness.
  \textit{American Journal of Applied Psychology}, \textit{14}(3), 101--112.
  \url{https://doi.org/10.11648/j.ajap.20251403.14}

\item Palos-Sánchez, P.~R., Baena-Luna, P., Badicu, A., \& Infante-Moro, J.~C. (2022).
  Artificial intelligence and human resources management~: A bibliometric analysis.
  \textit{Applied Artificial Intelligence}, \textit{36}(1), e2145631.
  \url{https://doi.org/10.1080/08839514.2022.2145631}

\item Soulami, M., Benchekroun, S., \& Galiulina, A. (2024). Exploring how AI adoption
  in the workplace affects employees~: A bibliometric and systematic review.
  \textit{Frontiers in Artificial Intelligence}, \textit{7}, 1473872.
  \url{https://doi.org/10.3389/frai.2024.1473872}

\item Wagh, R.~K. (2025). Exploring the role of AI and automation in human resource
  training~: Opportunities and challenges. \textit{International Journal of Research
  in Human Resource Management~(IJRHRM)}, \textit{7}(2), 245--251.
  \url{https://doi.org/10.33545/26633213.2025.v7.i2c.349}

\end{enumerate}

\end{document}
